\chapter{Design and Implementation}
\label{chap:design-and-implementation}

The motivation and trade-offs for each technique are established in \cref{subsec:bg-refproc-tradeoffs}; this chapter describes their design and implementation. Three optimisations target \ac{ZGC}'s processing of \texttt{WeakReference} objects without a \texttt{ReferenceQueue}: skip-enqueue separation (\cref{sec:di-sep}), the dynamic array (\cref{sec:di-dyn}), and the specialised clear path (\cref{sec:di-clear}). In addition, an exploratory weak-field implementation (\cref{sec:di-wf}) was developed to evaluate an alternative way of expressing weak semantics that avoids the overhead of \texttt{WeakReference} objects entirely.

\Cref{tab:mechanism-summary} provides a compact overview of how the four mechanisms differ across the key design dimensions, meanwhile \Cref{tab:mechanism-cost-benefit} summarises the expected costs and benefits of each mechanism.

\begin{table}[t]
\centering
\caption[Summary comparison of the four mechanisms]{Comparison of the four mechanisms along key design dimensions. The representation column describes how weak semantics are expressed. The discovery structure column describes the data structure used during the concurrent mark phase to record discovered references. The enqueue path column describes whether discovered references are enqueued to the pending list for the \texttt{ReferenceHandler} thread. The clear path column describes the operations performed to clear the referent field of an unreachable reference during processing.}
\label{tab:mechanism-summary}
\small
\begin{tabularx}{\textwidth}{@{}>{\hsize=1.0\hsize\raggedright\arraybackslash}X
                              >{\hsize=1.0\hsize\raggedright\arraybackslash}X
                              >{\hsize=1.0\hsize\raggedright\arraybackslash}X
                              >{\hsize=1.0\hsize\raggedright\arraybackslash}X
                              >{\hsize=1.0\hsize\raggedright\arraybackslash}X@{}}
\toprule
\textbf{Mechanism} & \textbf{Representation} & \textbf{Discovery structure} & \textbf{Enqueue path} & \textbf{Clear path} \\
\midrule
Baseline (\texttt{none})
  & Independent objects
  & Intrusive linked list
  & Pending list
  & Load barrier + virtual call + \ac{CAS} \\
\addlinespace
Skip-enqueue sep. (\texttt{sep})
  & Independent objects
  & Separate linked lists per list type
  & Queue-less refs bypass pending list
  & Load barrier + virtual call + \ac{CAS} \\
\addlinespace
Dynamic array (\texttt{dyn})
  & Independent objects
  & Contiguous array
  & Pending list
  & Load barrier + virtual call + \ac{CAS} \\
\addlinespace
Clear path (\texttt{clear\_path})
  & Independent objects
  & Intrusive linked list
  & Pending list
  & Pre-loaded data + direct call + plain store \\
\addlinespace
Weak fields (\texttt{weak\_fields})
  & Annotated object fields
  & Contiguous array per \ac{GC} worker
  & No pending list
  & Pre-loaded data + direct call + \ac{CAS}\\
\bottomrule
\end{tabularx}
\normalsize
\end{table}

\begin{table}[t]
\centering
\caption[Cost and benefit comparison of the four mechanisms]{Comparison of the expected costs and benefits of the four mechanisms. The discovery structure column describes the data structure used during the concurrent mark phase to record discovered references. ``Clear path'' describes how the referent field is cleared once the reference is found to be unreachable.}
\label{tab:mechanism-cost-benefit}
\small
\begin{tabularx}{\textwidth}{@{}>{\hsize=1.0\hsize\raggedright\arraybackslash}X
                              >{\hsize=1.0\hsize\raggedright\arraybackslash}X
                              >{\hsize=1.0\hsize\raggedright\arraybackslash}X@{}}
\toprule
\textbf{Mechanism} & \textbf{Expected benefit} & \textbf{Expected cost} \\
\midrule
Skip-enqueue sep. (\texttt{sep})
  & Shorter reference-processing time and less work for the \texttt{ReferenceHandler} thread for queue-less weak references
  & Increased code complexity and slightly larger code objects\\
\addlinespace
Dynamic array (\texttt{dyn})
  & Shorter reference-processing time due to better cache locality and pre-loaded field data
  & Higher memory footprint and longer discovery time due to dynamic array management\\
\addlinespace
Clear path (\texttt{clear\_path})
  & Shorter reference-processing time due to fewer and simpler operations for clearing the referent field
  & Increased code complexity and potential for subtle bugs in future development\\
\addlinespace
Weak fields (\texttt{weak\_fields})
  & Shorter total \ac{GC} cycle time and lower memory footprint due to avoiding \texttt{WeakReference} objects entirely
  & Immensely increased code complexity and an API-level separation between weak semantics with and without callbacks\\
\bottomrule
\end{tabularx}
\normalsize
\end{table}

\section{Skipping Enqueue Path}
\label{sec:di-sep}
As established in \cref{subsec:bg-refproc-tradeoffs}, skip-enqueue separation eliminates pending-list work for queue-less weak references by routing them to a separate discovered list that bypasses the enqueue stage entirely.

\subsection{Implementation of Skipping Enqueue Path}
\label{subsec:di-sep-impl}
In order to skip enqueueing weak references without a \texttt{ReferenceQueue} the \texttt{ZReferenceProcessor} class was modified to check if the queue field of a weak reference being discovered is equal to the sentinel value \texttt{ReferenceQueue.NULL\_QUEUE} before adding it to the discovered list. If the queue field is equal to \texttt{ReferenceQueue.NULL\_QUEUE} (the weak reference was not created with a queue), it is added to a separate discovered list for weak references without a \texttt{ReferenceQueue} instead of being added to the regular discovered list. This separate discovered list is processed independently by the \ac{GC} threads and none of its weak references are enqueued to the pending list for the \texttt{ReferenceHandler} thread. Each \ac{GC} thread maintains its own separate discovered list for weak references without a \texttt{ReferenceQueue} to avoid contention between threads. See \cref{lst:null-queue-fast-path} for a pseudocode summary and \cref{app:source-refproc} for the actual implementation.

The discovered lists are separated rather than handled through conditional branches in a single processing loop so that two specialised reference-processing functions can be used: one for queue-backed references, which includes enqueue handling, and one for queue-less references, which omits it. This allows the queue-less path to avoid unnecessary per-reference checks during processing and keeps the no-enqueue path structurally simpler.

\begin{listing}[t]
\begin{lstlisting}[language=C++]
bool discover_reference(oop ref_obj, ReferenceType type) {
    zaddress ref = to_zaddress(ref_obj), // ZGC-coloured pointer

    if (!should_discover(ref, type)) {   // skip if already discovered or ineligible
        return false,
    }

    if (type == REF_WEAK && !has_reference_queue(ref)) {
        // Queue-less weak ref: bypasses pending list and
        // ReferenceHandler thread entirely
        weak_no_queue_list_per_worker.append(ref),
    } else {
        // All other types (and queue-backed weak refs) follow
        // the standard discovered-list and enqueue pipeline
        discovered_list_per_worker.append(ref),
    }

    mark_discovered(ref), // set discovered field to prevent re-discovery
    return true,
}
\end{lstlisting}
\caption[Queue-less weak discovery split]{Pseudocode of the null-queue fast path during discovery. Queue-less weak references are routed to a separate discovered list and processed by a no-enqueue pipeline (see \cref{app:source-refproc}).}
\label{lst:null-queue-fast-path}
\end{listing}

\section{Dynamic Array for Discovered References}
\label{sec:di-dyn}
As established in \cref{subsec:bg-refproc-tradeoffs}, replacing the intrusive linked list with a contiguous dynamic array improves spatial locality during the processing traversal~\cite{drepper07-memory,jones11-gc-handbook}. The array is used for weak references without a \texttt{ReferenceQueue}, and it also enables pre-loading field data during discovery.

\subsection{Implementation of Dynamic Array for Discovered References}
\label{subsec:di-dyn-impl}

The implemented dynamic array, called \texttt{ZWeakRefArray}, stores discovered references in a contiguous layout rather than linking them via the \texttt{discovered} field. The array is allocated on the C heap and grows dynamically as needed.

Each \ac{GC} worker thread maintains its own \texttt{ZWeakRefArray} instance to avoid contention. During discovery, when a reference should be added to a discovery list, its data is appended to the worker's array instead of being appended to a linked list. The reference's \texttt{discovered} field is set to a self-referencing pointer to mark the reference as discovered, preventing duplicate discovery.

During processing, the array is iterated sequentially using index-based access. Because the elements are stored contiguously in memory, this iteration improves cache line utilisation compared to following pointers scattered across the heap~\cite{drepper07-memory}. After processing, the array is cleared for the next \ac{GC} cycle but retains enough capacity to at least hold the same number of references that were processed but not made inactive. This avoids unnecessary reallocations in subsequent cycles when the number of discovered references is relatively stable.

See \cref{lst:dynamic-array-processing} for a pseudocode summary and \cref{app:source-dyn,app:source-refproc} for the actual implementation.

\begin{listing}[t]
\begin{lstlisting}[language=C++]
void discover_reference(oop ref_obj, ReferenceType type) {
    //...
    ZWeakRefData data,
    data.set_from_ref(ref),    // pre-load referent address and value
    worker_array.append(data), // O(1) amortised; array grows if needed
    //...
}

void process_queue_less_weak_refs(ZWeakRefArray& arr) {
    size_t dropped = 0,
    for (size_t i = 0, i < arr.length(), i++) {
        const ZWeakRefData& data = arr.at(i), // sequential: cache-friendly
        data.mark_not_discovered(),            // clear discovered field first

        if (try_make_inactive(data)) { // clear referent if unreachable
            cleared_weak_no_queue_count++,
        } else {
            dropped++, // reference still active; keep slot
        }
    }
    arr.clear_and_reserve(dropped), // retain capacity for next cycle
}
\end{lstlisting}
\caption[Dynamic-array weak processing]{Pseudocode for the \texttt{ZWeakRefArray}-backed queue-less weak-reference path. Discovery appends pre-loaded fields and processing scans the array sequentially (see \cref{app:source-dyn,app:source-refproc}).}
\label{lst:dynamic-array-processing}
\end{listing}

\section{Optimised Clear Path}
\label{sec:di-clear}

As established in \cref{subsec:bg-refproc-tradeoffs}, the standard \ac{ZGC} clear path performs three operations per reference that can be simplified or eliminated for queue-less old-generation \texttt{WeakReference} objects~\cite{openjdk-zreferenceprocessor-cpp,openjdk-zbarrier-cpp}:
\begin{enumerate}
    \item A load barrier to read the referent from the reference object.
    \item A virtual call to determine the reference type (soft, weak, final, or phantom).
    \item A \ac{CAS} operation via a \ac{ZGC} barrier to atomically set the referent field to a coloured null pointer.
\end{enumerate}

For weak references without a \texttt{ReferenceQueue}, each of these operations can be simplified or eliminated:
\begin{enumerate}
    \item When combined with the dynamic array optimisation, the referent field address and its value can be pre-loaded during discovery and stored in the array entry. This avoids redundant memory loads during processing. Pre-loading is safe because the only concurrent application-level operations on a \texttt{Reference}'s referent are clearing and enqueueing, both of which set the referent field to null rather than to a new non-null value~\cite{openjdk-reference-java,openjdk-zreferenceprocessor-cpp}.
    \item The reference type is known to be \texttt{REF\_WEAK} at the call site, eliminating the need for a virtual call to determine the type.
    \item The only concurrent modification to the referent field is an application thread setting it to null~\cite{openjdk-reference-java,openjdk-zbarrier-cpp,openjdk-zreferenceprocessor-cpp}. This means no thread can set the field to a new non-null value, which means that the \ac{CAS} operation for clearing the referent can be replaced with a simple store of a coloured null pointer.
\end{enumerate}

These simplifications reduce both instruction count and memory-access overhead for each weak reference without a \texttt{ReferenceQueue} during processing, which should in turn reduce reference-processing time and \ac{GC} cycle time.

\subsection{Implementation of Optimised Clear Path}
\label{subsec:di-clear-impl}

The implementation introduces a new function for determining whether a queue-less weak reference's referent should be cleared (see \cref{app:source-refproc}). This function takes a data struct (either from the dynamic array or constructed on the fly from the linked list) containing the pre-loaded field addresses and values. This new function has similar semantics to the standard function and the \ac{ZGC} barriers it uses, but it is specialised for the queue-less weak reference case and can therefore use more efficient operations. See \cref{sec:bg-zgc} for more details on \ac{ZGC} barriers.

The new function checks whether the referent is strongly live via metadata in \ac{ZGC}'s heap management. If the referent is not strongly live and resides in the old generation, its referent field is cleared using a direct pointer store rather than a \ac{CAS}. This is safe under the constrained update model described above. If the referent is still live, either because it is strongly marked or because it resides in the young generation, the referent field is recoloured with the correct pointer metadata bits. That recolouring has to use \ac{CAS} so that a null value is not overwritten with a non-null value. Additionally, if the referent is young and the young generation is currently in its mark phase, the referent is marked to preserve cross-generational liveness~\cite{openjdk-zreferenceprocessor-cpp,openjdk-zbarrier-cpp,openjdk-zgeneration-cpp}.

The \texttt{discovered} field is also cleared using a direct store to the field address rather than through an atomic access function, since no other thread modifies this field after discovery in this processing path~\cite{openjdk-zreferenceprocessor-cpp}.

See \cref{lst:gen-opt-fast-inactive} for a pseudocode summary and \cref{app:source-refproc} for the actual implementation.

\begin{listing}[t]
\begin{lstlisting}[language=C++]
bool try_make_inactive_fast(const ZWeakRefData& data) {
    // Array variant: referent address pre-loaded at discovery (no barrier needed)
    zaddress referent = data.referent_addr,

    // Linked-list variant: must load through a barrier to get the live address
    zaddress referent = load_barrier(data.referent_field_addr),

    // Young referents cannot be cleared here: they are kept alive by
    // generational rules and are only processed in a young-gen collection
    if (is_young(referent)) {
        if (young_mark_is_active()) {
            mark_young_root(referent), // preserve cross-generational liveness
        }
        return false, // keep active
    }

    // Old referent is dead: no thread can store a non-null value back, so
    // a plain store is safe (no CAS required)
    if (!is_strongly_live(referent)) {
        *data.referent_field_addr = to_zpointer(zaddress::null),
        return true, // became inactive
    }

    // Referent is still live: update pointer colour bits for the current
    // GC phase, but avoid overwriting a null written by an application thread
    zpointer observed = data.referent_field_value,
    zpointer recoloured = remap_with_good_metadata(referent),
    cas_referent_if_non_null(data.referent_field_addr, observed, recoloured),
    return false,
}
\end{lstlisting}
\caption[Optimised-clear-path fast inactive check]{Pseudocode for the specialised queue-less weak-reference inactive check. The old-generation clear path uses direct stores, while recolouring uses \ac{CAS} to avoid null-to-non-null races (see \cref{app:source-refproc}).}
\label{lst:gen-opt-fast-inactive}
\end{listing}

\section{Weak Fields}
\label{sec:di-wf}

As established in \cref{subsec:bg-refproc-tradeoffs}, replacing \texttt{WeakReference} wrapper objects with annotated fields should reduce both reference-processing overhead and heap allocation pressure. To answer \cref{rq:weak-fields}, an exploratory weak-field mechanism was implemented. The mechanism allows developers to annotate object fields as weak, indicating that the \ac{GCr} should treat those fields with weak semantics. The \ac{GCr} then applies weak semantic rules directly to the annotated field without requiring a separate \texttt{WeakReference} object.

\subsection{Implementation of Weak Fields}
\label{subsec:di-wf-impl}

The implementation includes changes in three layers of the \ac{JVM}: a Java-level annotation, \ac{GCr}-side field discovery and processing, and compiler/interpreter support for the annotation metadata.

\subsubsection{Annotation}

A new runtime-retained field annotation, \texttt{@java.lang.ref.weak}, was added to the \texttt{java.lang.ref} package. The annotation targets only fields and is retained at runtime so that the VM can inspect it during class loading. See \cref{lst:weak-annotation} for the annotation definition. During class-file parsing, the VM recognises the \texttt{@weak} annotation and sets a \texttt{weak} flag in the field's metadata.

\begin{listing}[t]
\begin{lstlisting}[language=Java]
package java.lang.ref,

@Retention(RetentionPolicy.RUNTIME) // visible to the VM at class-loading time
@Target(ElementType.FIELD)          // applies to fields only, not classes or methods
public @interface weak {}
\end{lstlisting}
\caption[Weak-field annotation]{The \texttt{@weak} annotation. Applied to a field, it instructs the \ac{GCr} to treat the field value with weak-reference semantics (see \cref{app:source-wf-annotation}).}
\label{lst:weak-annotation}
\end{listing}

\subsubsection{Discovery}

Weak-field discovery occurs during \ac{ZGC}'s old-generation marking phase. As the marking closure iterates over each object's reference fields, it checks whether the current field carries the \texttt{weak} flag in its metadata. If the field is annotated as weak, the marking closure diverts it to the weak-field discovery path instead of applying the normal strong mark barrier~\cite{openjdk-zreferenceprocessor-cpp}.

The discovery function first checks whether the field should be discovered at all. Fields that already point to null or to a strongly live object (identified by their \ac{ZGC} pointer metadata bits) are skipped, since no clearing action will be needed for them. Fields in young-generation objects are also skipped, because young-generation weak-field processing is not supported by this implementation. If the field passes these checks, its address and current value are recorded as a \texttt{ZWeakFieldData} struct and appended to the \ac{GC} worker thread's per-worker \texttt{ZWeakFieldArray}. See \cref{lst:weak-field-discovery} for a pseudocode summary and \cref{app:source-wf-refproc} for the actual implementation.

Each \ac{GC} worker thread maintains its own \texttt{ZWeakFieldArray} instance to avoid contention. The array is structurally identical to the \texttt{ZWeakRefArray} used for queue-less weak references: it stores entries contiguously on the C heap and grows dynamically (see \cref{app:source-wf-array}).

\begin{listing}[t]
\begin{lstlisting}[language=C++]
// Called by the marking closure for every reference field of every live object
void do_oop(oop* p) {
    if (is_weak_annotated_field(p)) {  // check weak flag in field metadata
        if (discover_weak_field(p)) {
            return, // diverted to weak path, skip strong marking
        }
    }
    // ... normal strong marking ...
}

bool discover_weak_field(oop* field) {
    encountered_weak_fields_count++,

    if (is_young(field)) {
        return false, // young-generation weak fields are not supported
    }

    if (is_null(field) || is_strongly_live(field)) {
        return false, // no clearing needed: already null or strongly reachable
    }

    // Record field address and current pointer value for processing
    discovered_weak_fields_count++,
    ZWeakFieldData data(field),
    worker_weak_field_array.append(data), // per-worker; avoids contention
    return true,
}
\end{lstlisting}
\caption[Weak-field discovery]{Pseudocode for weak-field discovery during marking. Annotated fields are diverted from the strong mark barrier and recorded in a per-worker array (see \cref{app:source-wf-refproc}).}
\label{lst:weak-field-discovery}
\end{listing}

\subsubsection{Processing}

During the reference-processing phase, each \ac{GC} worker iterates through its \texttt{ZWeakFieldArray} and applies a clearing function to each entry. The function uses the data recorded during discovery. If the field still contains the discovered value and the referent is no longer strongly live, the routine clears the field by storing a coloured null pointer. If the field still contains the discovered value but the referent remains live, it recolours the pointer with the correct metadata bits. If the field value has changed since discovery, no action is required, since only a strongly reachable referent could have been stored in the field after discovery. Apart from this extra check for post-discovery field updates, the function follows the same overall strategy as the optimised clear path described in \cref{sec:di-clear}. After processing, the array is cleared but retains enough capacity to hold the entries that were not cleared, following the same capacity-retention strategy as \texttt{ZWeakRefArray}. See \cref{lst:weak-field-processing} for a pseudocode summary and \cref{app:source-wf-refproc} for the actual implementation.

\begin{listing}[t]
\begin{lstlisting}[language=C++]
void process_worker_discovered_weak_fields(ZWeakFieldArray& arr) {
    size_t dropped = 0,
    for (size_t i = 0, i < arr.length(), i++) {
        const ZWeakFieldData& data = arr.at(i), // sequential scan
        if (try_make_inactive_fast(data)) { // clear dead or recolour live
            cleared_weak_fields_count++,
        } else {
            dropped++, // field still active; keep slot for next cycle
        }
    }
    arr.clear_and_reserve(dropped), // retain capacity for next cycle
}
\end{lstlisting}
\caption[Weak-field processing]{Pseudocode for weak-field processing. Each entry is passed to the clear barrier, which clears dead referents and recolours live ones (see \cref{app:source-wf-refproc}).}
\label{lst:weak-field-processing}
\end{listing}

\subsubsection{Compiler and Interpreter Support}
The \texttt{@weak} annotation metadata is also made available to the \ac{JIT} compilers and interpreter. This is necessary to ensure that any loads or stores to a weak field outside \ac{GC} processing are handled correctly with respect to weak semantics. The implementation adds a decorator to loads and stores that target a field carrying the \texttt{weak} flag in its metadata. \ac{ZGC} then inserts barriers that apply the required semantics to all accesses with this decorator. Without the load decorator, a load from a weak field could return a non-null reference whose referent will be cleared by the \ac{GCr}. The store decorator is needed so that the \ac{GCr} does not record the previous value of the field, unlike in the default strong-reference case used to preserve the \ac{SATB} marking semantics. See \cref{subsec:bg-zgc-decorators} for more details on decorators and \cref{app:source-wf-compiler} for the implementations.
